\documentclass[12pt]{article}
\usepackage[T1]{fontenc}
\usepackage[utf8]{inputenc}
\usepackage{lmodern}
\usepackage{textcomp}
\usepackage{enumitem}
\usepackage{geometry}
\usepackage{graphicx}
\usepackage{amsmath, amssymb}
\geometry{margin=1in}
\begin{document}
\begin{center}
Chemistry - Graduate Level Assessment 19 \
Total Marks: 40 \
Answer all questions.
\end{center}

\section*{Multiple Choice (10 marks)}
\begin{enumerate}[label=\arabic*.]
\item The rotational spectrum of HF has lines 41.9 cm^-1 apart. Calculate the moment of inertia and the bond length of this molecule.
\begin{enumerate}[label=(\alph*)]
    \item 1.256 \ensuremath{\times} 10^-46 kg $m^2$ and 9.100 \ensuremath{\times} 10^-2 nm
    \item 1.336 \ensuremath{\times} 10^-47 kg $m^2$ and 8.250 \ensuremath{\times} 10^-2 nm
    \item 1.336 \ensuremath{\times} 10^-46 kg $m^2$ and 9.196 \ensuremath{\times} 10^-1 nm
    \item 1.410 \ensuremath{\times} 10^-47 kg $m^2$ and 9.250 \ensuremath{\times} 10^-2 nm
    \item 1.256 \ensuremath{\times} 10^-47 kg $m^2$ and 8.196 \ensuremath{\times} 10^-2 nm
\end{enumerate}
\item $ \text {Calculate the energy of a photon for a wavelength of } 100 \mathrm{pm} \text { (about one atomic diameter). } $\$
\begin{enumerate}[label=(\alph*)]
    \item 4 x 10^-15 J
    \item 1.2 x 10^-15 J
    \item 5 x 10^-15 J
    \item 3.5 x 10^-15 J
    \item 2 $10^{-15} \mathrm{~J}$
\end{enumerate}
\item Less than 1/1000 of the mass of any atom is contributed by
\begin{enumerate}[label=(\alph*)]
    \item the protons, electrons and neutrons
    \item the electrons and protons
    \item the nucleus
    \item the protons
    \item the neutrons
\end{enumerate}
\item What is the electromotive force (emf) of a system when K, the equilibrium constant for the reaction = Q, the concentration quotient?
\begin{enumerate}[label=(\alph*)]
    \item -2
    \item 1
    \item -0.5
    \item 0.5
    \item 2
\end{enumerate}
\item What is the heat capacity of uranium? Its atomic heat is 26 J/mole and its heat capacity is 238 g/mole.
\begin{enumerate}[label=(\alph*)]
    \item 0.11 J/g^\ensuremath{\circC}
    \item 0.20 J/g^\ensuremath{\circC}
    \item 0.13 J/g^\ensuremath{\circC}
    \item 0.18 J/g^\ensuremath{\circC}
    \item 0.30 J/g^\ensuremath{\circC}
\end{enumerate}
\end{enumerate}

\section*{True / False (10 marks)}
\textit{Answer True or False}
\begin{enumerate}[label=\arabic*.]
\setcounter{enumi}{5}
\item True or False: The oxidation number of chlorine in the perchlorate ion, ClO 4-, is +5.
\item True or False: The strongest base in liquid ammonia is ammonia itself, NH 3.
\item True or False: An atom of 52 Cr contains 28 protons, 24 neutrons, and 24 electrons.
\item True or False: The compounds formed from Sn and F, P and H, and Si and O are SnF\_4, PH\_3, and Si\_2 O\_5, respectively.
\item True or False: A T$_{2}$ value of 15 ms results in a linewidth of approximately 52.8 Hz in NMR spectroscopy.
\end{enumerate}

\section*{Long Form Response (20 marks)}
\textit{Answer each question with a clear and complete explanation.}
\begin{enumerate}[label=\arabic*.]
\setcounter{enumi}{10}
\item Discuss the reasons why the 1 H NMR spectrum of 12 CHCl 3 appears as a singlet, focusing on the implications of the electric quadrupole moments of 35 Cl and 37 Cl.
\item Given an unknown element X with an atomic weight of 210.197 amu and two isotopes, ^210 X and ^212 X with respective masses of 209.64 and 211.66 amu, calculate the relative abundance of each isotope and discuss the implications of your findings on the element's stability.
\end{enumerate}

\vfill
\noindent\textit{End of Paper}
\end{document}